\originalchapter{量子光场与非经典态}
\originalsection{场量子化、相干态与光子统计}
对选定模式, $H=\hbar\omega(a^\dagger a+1/2)$,$[a,a^\dagger]=1$. 一般空间场为模式和, 每项有电场幅度$\sqrt{\hbar\omega/(2\epsilon_0V)}$乘模式函数. 实测场需要模式、频段、偏振与探测时间的约定; 一个$a$不能表示所有空间点的独立光子.

相干态由$a\ket\alpha=\alpha\ket\alpha$定义. 展开$\sum_nc_n\ket n$得$c_{n+1}=\alpha c_n/\sqrt{n+1}$, 归一化后
\[\ket\alpha=\e^{-|\alpha|^2/2}\sum_{n=0}^{\infty}\frac{\alpha^n}{\sqrt{n!}}\ket n,\qquad P_n=\e^{-\bar n}\frac{\bar n^n}{n!},\quad\bar n=|\alpha|^2.\]
所以$\operatorname{Var}n=\bar n$. 相干态的相位相干来自不同$n$的确定相对相位; 对相位均匀平均后的Poisson混合态有相同光子数分布, 却有$\langle a\rangle=0$. 热态$P_n=\bar n^n/(1+\bar n)^{n+1}$, 方差$\bar n(1+\bar n)$; 数态方差零.

\originalsection{正交分量与压缩}
定义$X=(a+a^\dagger)/2,Y=(a-a^\dagger)/(2\ii)$, 则$[X,Y]=\ii/2$且$\Delta X\Delta Y\ge1/4$. 真空与相干态方差均为$1/4$. 对$S(r)=\exp[r(a^2-a^{\dagger2})/2]$, 用参数微分和对易子得$S^\dagger aS=a\cosh r-a^\dagger\sinh r$, 故$X\mapsto\e^{-r}X,Y\mapsto\e^rY$. 压缩真空的两个方差为$\e^{-2r}/4,\e^{2r}/4$, 平均光子数$\sinh^2r$.

\begin{center}\begin{tikzpicture}[scale=1.15,>=Stealth]
\draw[->](-2,0)--(2.4,0)node[right]{$X$};\draw[->](0,-1.8)--(0,1.8)node[above]{$Y$};
\draw[dashed](0,0)circle(1);\draw[thick,structurecolor](0,0)ellipse(.5 and 1.5);
\node at(1.45,1.3){真空};\node[structurecolor]at(-1.2,1.6){压缩态};
\end{tikzpicture}\end{center}
图为正交分量噪声的示意轮廓, 不是原子位置的椭圆; 面积不因压缩而减少. 参数放大器以泵浦能量产生光子对, 可用$H_I=\ii\hbar\kappa(a^{\dagger2}-a^2)/2$产生这种变换. 泵浦耗尽、损耗和相位抖动都会降低可见压缩.

\originalsection{同相探测与准概率}
把待测模式$a$和强局域振荡器$b\simeq\beta\e^{\ii\phi}$送入50:50分束器. 两输出$c=(a+b)/\sqrt2,d=(b-a)/\sqrt2$的光子数差为$c^\dagger c-d^\dagger d=a^\dagger b+b^\dagger a\simeq2|\beta|X_\phi$, $X_\phi=(a\e^{-\ii\phi}+a^\dagger\e^{\ii\phi})/2$. 改变局域相位即可测不同投影, 重建Wigner函数.

Wigner函数可为负, 因为它不是$X,Y$同时具有确定值的普通概率. 在本卷约定中真空$W(X,Y)=2\e^{-2(X^2+Y^2)}/\pi$, 相干态平移中心, 压缩态变成椭圆. 数态$n=1$的$W=(2/\pi)[4(X^2+Y^2)-1]\e^{-2(X^2+Y^2)}$, 原点为负. 不能把Wigner的负值与负的探测概率混淆: 任何一个正交分量的边缘分布仍为非负.

有限效率$\eta$等价于分束器混入真空, 方差变为$\eta V+(1-\eta)/4$. 即使入射噪声趋零, 输出也有$(1-\eta)/4$的下限. 实验中的“低于散粒噪声”必须使用相同增益与探测效率的真空基准.
